% ======================================================================
% Appendix: proofs and measurement details for the two-solution theory
% (Sections \ref{sec:tworoutes}-\ref{sec:whydamping}).  Assumes
% \appendix was already issued (sec7_appendix.tex).
% ======================================================================
\section{Proofs and Measurement Details for the Two-Solution Analysis}
\label{app:tworoutes}

\subsection{Deferred Proofs}

\begin{proof}[Proof of Lemma~\ref{lem:clones}]
Induction on $t$.  At initialization all messages are zero.  If all
corresponding messages are equal at $t-1$, the two copies receive equal
inputs and hold equal tables, so their $R^{t}$ are equal; the two $Q^{t}$
from $X_i$ toward the pair differ only by which copy's $R^{t-1}$ is
excluded, and those are equal by hypothesis.
\end{proof}

\begin{proof}[Proof of Lemma~\ref{lem:decode}]
Each copy forwards half of the same row:
$R^{t}_{F'_{ij}\to X_i}(v) = \tfrac12 C_{ij}(v, u^{*t}_j) +
\text{const}$, and likewise for $F''_{ij}$ with the same minimizer
(Lemma~\ref{lem:clones}).  Summing the two copies over all neighbors
reconstructs $\sum_j C_{ij}(v, u^{*t}_j)$ up to an additive constant,
which does not affect the argmin.
\end{proof}

\begin{proof}[Proof of Lemma~\ref{lem:ema}]
Solving $m_{\mathrm{even}} = \lambda m_{\mathrm{odd}} + (1-\lambda)u$ and
$m_{\mathrm{odd}} = \lambda m_{\mathrm{even}} + (1-\lambda)v$ gives
$m_{\mathrm{even}} = (u+\lambda v)/(1+\lambda)$ and
$m_{\mathrm{odd}} = (v+\lambda u)/(1+\lambda)$; the deviation from this
alternating pair contracts by $\lambda$ per step.
\end{proof}

\begin{proof}[Proof of Theorem~\ref{thm:lambdastar}]
Within the pattern the computed difference of each flipping message
alternates between the constants $\Delta_1, \Delta_2$, so by
Lemma~\ref{lem:ema} its damped difference converges to the alternating
pair $(\Delta_1 + \lambda\Delta_2)/(1+\lambda)$,
$(\Delta_2 + \lambda\Delta_1)/(1+\lambda)$.  Pattern consistency requires
the first in the upper regime and the second in the lower:
$\Delta_1 + \lambda\Delta_2 \ge (1+\lambda)\tau_U$ and
$\Delta_2 + \lambda\Delta_1 \le (1+\lambda)\tau_L$, i.e., the stated
bounds.  Under strict inequalities the pair sits strictly inside the
regimes, drives are locally constant, and each difference follows an
affine contraction of rate $\lambda$; hence the orbit exists and attracts
a neighborhood.  Fixed-point invariance:
$m = \lambda m + (1-\lambda)\widehat{m}$ iff $m = \widehat{m}$.
\end{proof}

\begin{proof}[Proof of Proposition~\ref{prop:kn}]
By symmetry of the instance and zero initialization (the induction of
Lemma~\ref{lem:clones} extended over the automorphisms of the complete
graph), all $Q$-differences are equal at every iteration; call the
common value $q_t$.  Each split table is $5$ on the diagonal and $0$ off
it, so its difference map is $-\operatorname{clip}(q, -5, 5)$ with
transition interval $[-5, 5]$, and a $Q$-message sums the unary gap and
the $R$-differences of $2(n-1) - 1 = 2n-3$ incoming copies, giving the
stated scalar recurrence.  While $|q_t| \ge 5$ the computed drive
alternates between $\Delta_1 = 0.01 + 5(2n-3)$ and
$\Delta_2 = 0.01 - 5(2n-3)$; Theorem~\ref{thm:lambdastar} with
$\tau_U = -\tau_L = 5$ yields the threshold
$(10(n-2) - 0.01)/(10(n-1) + 0.01)$ after simplification, and local
attraction with two-step contraction $\lambda^{2}$.  Reachability from
zero initialization: inside the strip the recurrence is affine with slope
$\lambda - (1-\lambda)(2n-3)$, which is $\le -1$ exactly when
$\lambda \le (n-2)/(n-1)$; below the threshold the interior fixed point
is repelling and the oscillating iterates leave the strip in finitely
many steps, after which the committed alternation attracts.  At
$\lambda = 0.95$ (for $n = 12$) the interior slope is $-0.10$ and the
message differences converge to the interior fixed point, whose decoded
beliefs are exactly tied --- an artifact of the perfectly uniform
instance --- so assignments there are determined by tie-breaking rather
than by a surviving alternation.
\end{proof}

\begin{proof}[Proof of Corollary~\ref{cor:bipartite}]
For an edge with $i \in A$, $j \in B$: $x'_i = x_i$, $x'_j = y_j$,
$y'_i = y_i$, $y'_j = x_j$, so the two cross terms
$C_{ij}(x_i, y_j) + C_{ij}(y_i, x_j)$ are literally the two intra terms
$C_{ij}(x'_i, x'_j) + C_{ij}(y'_i, y'_j)$; unary terms match because
each variable contributes both of its values once on either side.  A
single-variable move in $x'$ corresponds to a single-layer move in
$(x,y)$ --- in layer $x$ if the variable is in $A$, in layer $y$
otherwise --- so layer-wise local minimality of $\mathrm{cost}_2$
transfers to coordinatewise local minimality of $\mathrm{cost}$ for $x'$
and $y'$ separately.
\end{proof}

\subsection{Measurement Protocol}

All measurements were produced with a vectorized reimplementation of the
synchronous two-phase schedule, verified to produce assignment
trajectories identical to the repository engines
(\texttt{BPEngine}, \texttt{SplitEngine}, \texttt{DampingEngine},
\texttt{DampingSCFGEngine}) on shared instances, with per-iteration
message normalization (without which undamped raw messages lose
floating-point precision by iteration ${\sim}30$).  Families: random
Erd\H{o}s--R\'enyi with $|D_i| = 10$ and uniform integer costs in
$[100, 200)$ at densities $0.25$ (sparse) and $0.6$ (dense); random
binary ($|D_i| = 2$) at density $0.25$; 3-coloring at density $0.15$
with conflict cost $10$; tie-breaking unary preferences uniform in
$[0, 10^{-2})$.  \emph{Period census}: $700$ iterations; the period of
the selected-assignment sequence is the smallest $p \le 64$ that repeats
over the final $150$ iterations (``detected periodic tail''); $12$--$20$
seeds per cell over $n \in \{10, 20, 40, 80\}$.  \emph{Commitment}: the
fraction of messages with a stable active minimizer, at iteration $9$,
$49$, and averaged over the final $50$ iterations, recorded for split
and unsplit runs of the same instances.  \emph{Audited mechanism
subset}: $34$ instances ($10$ sparse, $10$ coloring, $8$ binary, $6$
dense), of which $32$ had period-2 tails; on each we checked the
selection-rule closure on the limit pair (both directions, exact), the
selection-rule prediction on every seventh transition of the recorded
$100$-iteration tail, and layer-wise local minimality of
$\mathrm{cost}_2$ under all single-variable moves.  \emph{Damping scans}:
$19$ instances ($17$ oscillating undamped), grid
$\lambda \in \{0.02, 0.05, 0.08, 0.12, 0.2, 0.3, 0.4, 0.5, 0.6, 0.7,
0.8, 0.9\}$, $2000$ iterations, ``stable'' meaning the assignment is
constant over the final $300$.  \emph{Bipartite verification}: $40$
random complete-bipartite instances ($4{+}4$ variables, $|D_i| = 3$,
uniform integer costs in $[0,100)$), of which $8$ oscillated; on all
$8$, the identity of Corollary~\ref{cor:bipartite} held exactly and both
rephasings were coordinatewise local minima of the original problem
(raw snapshot costs $532$--$1044$ versus rephased costs $474$--$560$).
Raw results and scripts accompany the code
(\texttt{analysis\_oscillation/}).

\subsection{Period Census}

\begin{table}[t]
\centering
\caption{Fraction of runs whose selected-assignment sequence has a
detected periodic tail of period 1 / period 2 / longer / none (window
$150$, $p \le 64$).  MS-SCFG uses the symmetric split, no damping.}
\label{tab:census}
\scriptsize
\begin{tabular}{llrrrrrrrr}
\toprule
& & \multicolumn{4}{c}{MS-SCFG} & \multicolumn{4}{c}{MS} \\
\cmidrule(lr){3-6}\cmidrule(lr){7-10}
family & $n$ & P1 & P2 & P$>$2 & none & P1 & P2 & P$>$2 & none \\
\midrule
sparse dom-10 & 10 & .35 & .65 & 0 & 0 & .65 & 0 & .15 & .20 \\
 & 20 & 0 & 1.00 & 0 & 0 & .10 & .05 & 0 & .85 \\
 & 40 & 0 & 1.00 & 0 & 0 & 0 & 0 & 0 & 1.00 \\
 & 80 & 0 & 1.00 & 0 & 0 & 0 & 0 & 0 & 1.00 \\
dense dom-10 & 10 & .08 & .83 & .08 & 0 & .17 & .17 & 0 & .67 \\
 & 20 & .17 & .83 & 0 & 0 & .08 & 0 & 0 & .92 \\
 & 40 & 0 & 1.00 & 0 & 0 & 0 & 0 & 0 & 1.00 \\
coloring & 10 & .55 & .45 & 0 & 0 & .90 & .05 & .05 & 0 \\
 & 20 & .15 & .85 & 0 & 0 & .45 & .15 & .25 & .15 \\
 & 40 & 0 & 1.00 & 0 & 0 & 0 & .95 & 0 & .05 \\
 & 80 & 0 & 1.00 & 0 & 0 & 0 & 1.00 & 0 & 0 \\
binary & 10 & .85 & .15 & 0 & 0 & .80 & 0 & .20 & 0 \\
 & 20 & .95 & .05 & 0 & 0 & .95 & 0 & 0 & .05 \\
 & 40 & .80 & .15 & .05 & 0 & .65 & .15 & 0 & .20 \\
 & 80 & .35 & .65 & 0 & 0 & .20 & .20 & .10 & .50 \\
\bottomrule
\end{tabular}
\end{table}

Table~\ref{tab:census} gives the full census.  Pooled over all cells,
$191$ of $276$ split runs ($69\%$) had a detected period-2 tail, $83$
period-1, and two longer periods ($16$, $42$); the split runs never
wandered.  Restricted to the $n \ge 40$ cells of the domain-10 and
coloring families the period-2 count is $92/92$.  Unsplit coloring at
$n \ge 40$ also locks into period 2 --- its bounded conflict penalty
lets high-degree fields commit messages without splitting, consistent
with the theory (whenever commitment happens, the period is $1$ or
$2$); unsplit domain-10 runs almost always show no detected period.

\subsection{Damping Grid Audit}

Per-instance smallest grid $\lambda$ from which all larger tested values
were stable: random sparse $\{.12, .2, .2, .2, .2\}$; random dense
$\{.2, .2, .4\}$; coloring $\{.02, .3, .4, .4, .4\}$; binary
$\{.02, .08, .12, .2\}$.  Two dense instances were non-monotone (stable
at one grid value, oscillating at a larger one, stable thereafter),
which is why the ``all larger values'' criterion is used.  On instances
where the number of flipping variables was unchanged from the undamped
run --- a proxy for an unchanged pattern, not a verification of it ---
the residual swing of the damped $Q$-differences followed the
$(1-\lambda)/(1+\lambda)$ law of Lemma~\ref{lem:ema} within a mean
relative deviation of about $4\%$ (random families; coloring
re-equilibrates into different patterns at any tested $\lambda > 0$ and
the same-pattern premise fails there).  Greedy per-variable rephasing of
the limit pair (minimizing the sum of the two layers' ordinary costs)
improved that sum by, on average, $2427$ (sparse), $4096$ (dense),
$374$ (coloring), and $412$ (binary), flipping $6$--$36\%$ of the
variables.
